\section{石炭化の順序・原子比・時代と形成温度を比較する。}

\subsection{平成24年度 問26}
\label{q:h24-reg-applied-26}

\noindent\textbf{原典：} sources/H24/17\_applied\_science.pdf、PDF物理8頁。

\textbf{公式正答：} 4（single）。

\par\textbf{検討未確定：} 公式4を保持。200℃以下は典型的石炭化の代表説明であり、原典の古生代という年代だけの条件から全炭田の熱史を確定できない。無煙炭等の高温例も一次資料で確認。

\paragraph{問題の要旨} 石炭化の順序・原子比・時代と形成温度を比較する。

\Needspace{4\baselineskip}\paragraph{確認した条件・記述（独自要約）}\mbox{}\par

\noindent\textbf{原典確認：} S7\_050の29原典画像で34問の全五肢・否定・数値を照合。\par

\noindent\textbf{対象：} 石炭化の順序・原子比・時代と形成温度を比較する。\par

\paragraph{問われる核心} 石炭化の順序・原子比・時代と形成温度を比較する。

\paragraph{必要知識} 年代・ランク・植物・温度を別の軸で読む。日本の炭田年代を主にジュラ紀とする1、褐炭の位置を交換した2、H/C・O/Cを増加とする3、石炭紀の代表植生を針葉樹とする5は不適切。提示された相対比較では4を選ぶが、年代だけから石炭の最高到達温度を200℃以下と保証しない。

\paragraph{考え方・正答を選ぶ理由（査読待ち初稿）} 年代・ランク・植物・温度を別の軸で読む。日本の炭田年代を主にジュラ紀とする1、褐炭の位置を交換した2、H/C・O/Cを増加とする3、石炭紀の代表植生を針葉樹とする5は不適切。提示された相対比較では4を選ぶが、年代だけから石炭の最高到達温度を200℃以下と保証しない。

\Needspace{4\baselineskip}\paragraph{選択肢の検討（査読待ち初稿）}\mbox{}\par

\noindent\textbf{1：不適切} 日本の主な石炭はジュラ系。日本の代表的炭田には古第三系等が多い。少数の古い産出例を日本全体の主流としない。\par

\noindent\textbf{2：不適切} 亜瀝青・瀝青・褐炭の順で石炭化。基本の低ランクからの順は褐炭→亜瀝青炭→瀝青炭→無煙炭で、褐炭の位置が逆。\par

\noindent\textbf{3：不適切} 石炭化でH/C・O/Cが増加。脱水・脱炭酸・炭素濃集等で一般にはこれら比は減少する方向へ進む。\par

\noindent\textbf{4：条件・解釈を留保} 古生代石炭は一般に200℃以下。石炭化は埋没熱史による。USGSの概説には150–200℃を超えない典型の説明があるが、高ランク・熱変成では200–300℃の例もある。年代だけで上限は定めない。\par

\noindent\textbf{5：不適切} 石炭紀の主起源を針葉樹。石炭紀の代表的な湿地森林には大型のシダ植物等が多く、針葉樹だけを代表起源としない。\par

\paragraph{再利用する考え方} 地質環境・生成機構・物質の名称を対応させ、平均値と全称、典型例と例外、当時の分類と現在の知見を分ける。

\paragraph{共通理論} \hyperref[theory:ore:coal]{第\ref*{theory:ore:coal}節}

\par 参考：\cite{ore50_coal_ref,ore50_coal_temp_ref,ore50_japan_coal_ref}。

\par 難易度 4（校準初稿）、優先度未評価。

\section{地熱資源の利用域・貯留層・発電方式を比較する。}

\subsection{平成23年度 問26}
\label{q:h23-reg-applied-26}

\noindent\textbf{原典：} sources/H23/17\_applied\_science.pdf、PDF物理8頁。

\textbf{公式正答：} 2（single）。

\par\textbf{検討未確定：} 地温勾配と150℃区分は原典の代表値。全地域を固定した境界で分けず、具体地点の勾配・貯留温度と地化学平衡条件は追加照合対象。

\paragraph{問題の要旨} 地熱資源の利用域・貯留層・発電方式を比較する。

\Needspace{4\baselineskip}\paragraph{確認した条件・記述（独自要約）}\mbox{}\par

\noindent\textbf{原典確認：} S7\_050の29原典画像で34問の全五肢・否定・数値を照合。\par

\noindent\textbf{対象：} 地熱資源の利用域・貯留層・発電方式を比較する。\par

\paragraph{問われる核心} 地熱資源の利用域・貯留層・発電方式を比較する。

\paragraph{必要知識} 地熱利用には高温蒸気による発電だけでなく、中低温熱水の直接利用やバイナリー発電がある。非火山地域の温水利用もあるため火山域への限定2を選ぶ。流体貯留と循環には割れ目等の透水性が重要。地化学温度計は平衡・混合等の仮定を確認して地下温度を推定する。

\paragraph{考え方・正答を選ぶ理由（査読待ち初稿）} 地熱利用には高温蒸気による発電だけでなく、中低温熱水の直接利用やバイナリー発電がある。非火山地域の温水利用もあるため火山域への限定2を選ぶ。流体貯留と循環には割れ目等の透水性が重要。地化学温度計は平衡・混合等の仮定を確認して地下温度を推定する。

\Needspace{4\baselineskip}\paragraph{選択肢の検討（査読待ち初稿）}\mbox{}\par

\noindent\textbf{1：代表的条件で適切} 蒸気方式とバイナリー。蒸気でタービンを回す方式と、熱交換で低沸点の二次媒体を蒸発させる方式を区別する。\par

\noindent\textbf{2：不適切} 中・低温熱水の利用範囲。温室・暖房・温泉等への直接利用は火山周辺だけに限定されない。\par

\noindent\textbf{3：条件・解釈を留保} 広域と局所の地温勾配。放射性壊変・深部熱供給・火山等の熱源を考える。20–30、50–100℃/km以上は代表範囲で全地点の一定値ではない。\par

\noindent\textbf{4：代表的条件で適切} 割れ目に富む貯留層。流体が集まり循環するために有効な透水経路が必要で、亀裂や断層が重要となる。\par

\noindent\textbf{5：代表的条件で適切} 化学成分による温度推定。溶解平衡等の温度依存を利用する地化学温度計。冷却、混合、再平衡の影響は別に検討する。\par

\paragraph{再利用する考え方} 地質環境・生成機構・物質の名称を対応させ、平均値と全称、典型例と例外、当時の分類と現在の知見を分ける。

\paragraph{共通理論} \hyperref[theory:ore:geothermal]{第\ref*{theory:ore:geothermal}節}

\par 参考：\cite{ore50_geothermal_ref,ore50_reservoir_ref}。

\paragraph{関連問題} \hyperref[q:r03-reg-applied-20]{R03-REG-APPLIED-20}

\par 難易度 4（校準初稿）、優先度未評価。

\subsection{令和3年度 問20}
\label{q:r03-reg-applied-20}

\noindent\textbf{原典：} sources/R03/17\_applied\_science.pdf、PDF物理7頁。

\textbf{公式正答：} 4（single）。

\par\textbf{検討未確定：} 地温勾配の代表範囲と地化学温度計の平衡仮定は地域・化学種の条件を要する。原典の固定範囲を地球全域の厳密定数としない。

\paragraph{問題の要旨} 地熱貯留層の透水経路と温度利用を比較する。

\Needspace{4\baselineskip}\paragraph{確認した条件・記述（独自要約）}\mbox{}\par

\noindent\textbf{原典確認：} S7\_050の29原典画像で34問の全五肢・否定・数値を照合。\par

\noindent\textbf{対象：} 地熱貯留層の透水経路と温度利用を比較する。\par

\paragraph{問われる核心} 地熱貯留層の透水経路と温度利用を比較する。

\paragraph{必要知識} 発電・直接利用の方式と、地下の熱・流体・透水性を分ける。通常の熱水貯留層は流体が移動する割れ目等を持ち、乏しいとした4を選ぶ。高温だが透水性不足の岩体の利用には別の技術と条件が必要で、通常の貯留層の定義と混同しない。

\paragraph{考え方・正答を選ぶ理由（査読待ち初稿）} 発電・直接利用の方式と、地下の熱・流体・透水性を分ける。通常の熱水貯留層は流体が移動する割れ目等を持ち、乏しいとした4を選ぶ。高温だが透水性不足の岩体の利用には別の技術と条件が必要で、通常の貯留層の定義と混同しない。

\Needspace{4\baselineskip}\paragraph{選択肢の検討（査読待ち初稿）}\mbox{}\par

\noindent\textbf{1：条件・解釈を留保} 広域と局所の地温勾配。放射性壊変・深部熱供給・火山等の熱源を考える。20–30、50–100℃/km以上は代表範囲で全地点の一定値ではない。\par

\noindent\textbf{2：代表的条件で適切} 非火山地域も低温利用の対象。地殻の広域熱等による熱水も利用でき、火山周辺だけに限定されない。\par

\noindent\textbf{3：代表的条件で適切} 化学成分による温度推定。溶解平衡等の温度依存を利用する地化学温度計。冷却、混合、再平衡の影響は別に検討する。\par

\noindent\textbf{4：不適切} 割れ目に乏しい地熱貯留層。通常の開発可能な熱水循環には割れ目等の透水経路が重要。原典H23の富んだを乏しいへ反転したもの。\par

\noindent\textbf{5：代表的条件で適切} 蒸気方式とバイナリー。蒸気でタービンを回す方式と、熱交換で低沸点の二次媒体を蒸発させる方式を区別する。\par

\paragraph{再利用する考え方} 地質環境・生成機構・物質の名称を対応させ、平均値と全称、典型例と例外、当時の分類と現在の知見を分ける。

\paragraph{共通理論} \hyperref[theory:ore:geothermal]{第\ref*{theory:ore:geothermal}節}

\par 参考：\cite{ore50_geothermal_ref,ore50_reservoir_ref}。

\paragraph{関連問題} \hyperref[q:h23-reg-applied-26]{H23-REG-APPLIED-26}

\par 難易度 4（校準初稿）、優先度未評価。

\section{黒鉱の鉱物・海底熱水成因と母岩を比較する。}

\subsection{平成26年度 問29}
\label{q:h26-reg-applied-29}

\noindent\textbf{原典：} sources/H26/17\_applied\_science.pdf、PDF物理12頁。

\textbf{公式正答：} 4（single）。

\par\textbf{検討未確定：} 原典の生成年代を西黒沢期に限る全称と全鉱物組合せの例外は追加照合。代表地域の記録を全黒鉱の唯一の産状へ一般化しない。

\paragraph{問題の要旨} 黒鉱の鉱物・海底熱水成因と母岩を比較する。

\Needspace{4\baselineskip}\paragraph{確認した条件・記述（独自要約）}\mbox{}\par

\noindent\textbf{原典確認：} S7\_050の29原典画像で34問の全五肢・否定・数値を照合。\par

\noindent\textbf{対象：} 黒鉱の鉱物・海底熱水成因と母岩を比較する。\par

\paragraph{問われる核心} 黒鉱の鉱物・海底熱水成因と母岩を比較する。

\paragraph{必要知識} 黒鉱は金属を含む海底熱水と硫化物の沈殿に関係し、単なる熱による再結晶という4が不適切。グリーンタフは変質した火山岩・火山砕屑岩の地域的呼称で、緑色片岩という変成岩と取り違えない。

\paragraph{考え方・正答を選ぶ理由（査読待ち初稿）} 黒鉱は金属を含む海底熱水と硫化物の沈殿に関係し、単なる熱による再結晶という4が不適切。グリーンタフは変質した火山岩・火山砕屑岩の地域的呼称で、緑色片岩という変成岩と取り違えない。

\Needspace{4\baselineskip}\paragraph{選択肢の検討（査読待ち初稿）}\mbox{}\par

\noindent\textbf{1：代表的条件で適切} ZnS・PbS・重晶石等の黒色塊状鉱。黒鉱の代表組合せは閃亜鉛鉱・方鉛鉱・黄鉄鉱・重晶石等。四面銅鉱等の副成分もある。\par

\noindent\textbf{2：代表的条件で適切} 珪長質海底火山に関係。火山性塊状硫化物鉱床の一型で、海底熱水活動による鉱化を考える。\par

\noindent\textbf{3：条件・解釈を留保} 日本のグリーンタフ・西黒沢期。東北日本の中新世グリーンタフ地域と秋田県北部は代表産地。限られるという厳密な年代全称は当時の分類と例外の照合が残る。\par

\noindent\textbf{4：不適切} 単なる火成岩周辺の再結晶。これは接触変成の要点で、金属を運ぶ熱水の沈殿による黒鉱成因を説明しない。\par

\noindent\textbf{5：代表的条件で適切} キプロス型よりPb・Znに富む。珪長質に関係する黒鉱型と苦鉄質海洋地殻に関係するキプロス型の代表的組成差。\par

\paragraph{再利用する考え方} 地質環境・生成機構・物質の名称を対応させ、平均値と全称、典型例と例外、当時の分類と現在の知見を分ける。

\paragraph{共通理論} \hyperref[theory:ore:kuroko]{第\ref*{theory:ore:kuroko}節}

\par 参考：\cite{ore50_kuroko_ref,ore50_backarc_ref}。

\par 難易度 4（校準初稿）、優先度未評価。

\subsection{令和5年度 問23}
\label{q:r05-reg-applied-23}

\noindent\textbf{原典：} sources/R05/17\_applied\_science.pdf、PDF物理8頁。

\textbf{公式正答：} 3（single）。

\par\textbf{検討未確定：} スカルン金属種を高温と中低温だけで固定する原典4の一般化は追加照合。全形態分類の完全排他性も主張しない。

\paragraph{問題の要旨} 黒鉱のグリーンタフと緑色片岩を区別する。

\Needspace{4\baselineskip}\paragraph{確認した条件・記述（独自要約）}\mbox{}\par

\noindent\textbf{原典確認：} S7\_050の29原典画像で34問の全五肢・否定・数値を照合。\par

\noindent\textbf{対象：} 黒鉱のグリーンタフと緑色片岩を区別する。\par

\paragraph{問われる核心} 黒鉱のグリーンタフと緑色片岩を区別する。

\paragraph{必要知識} 黒鉱の母岩をグリーンタフと緑色片岩で取り違えた3を選ぶ。熱水成因が共通でも、斑岩銅の分散した細脈、鉱脈の通路、スカルンの化学交代、海底の黒鉱は位置と形態が違う。温度と金属種の対応は酸化還元・母岩・流体組成の条件も確認する。

\paragraph{考え方・正答を選ぶ理由（査読待ち初稿）} 黒鉱の母岩をグリーンタフと緑色片岩で取り違えた3を選ぶ。熱水成因が共通でも、斑岩銅の分散した細脈、鉱脈の通路、スカルンの化学交代、海底の黒鉱は位置と形態が違う。温度と金属種の対応は酸化還元・母岩・流体組成の条件も確認する。

\Needspace{4\baselineskip}\paragraph{選択肢の検討（査読待ち初稿）}\mbox{}\par

\noindent\textbf{1：代表的条件で適切} 熱水鉱床の形態分類。鉱脈、塊状、スカルン、斑岩銅等の代表タイプを区別する。全分類が完全に排他的とはしない。\par

\noindent\textbf{2：代表的条件で適切} 半深成岩と斑岩銅。浅所の斑岩質貫入岩とその熱水系に関係する。\par

\noindent\textbf{3：不適切} 黒鉱が東北の緑色片岩中。黒鉱の代表は中新世グリーンタフの海底火山・火山砕屑岩地域。広域変成の緑色片岩と名称を交換している。\par

\noindent\textbf{4：条件・解釈を留保} 温度によるスカルン金属の傾向。高温段階にW等、後期・遠方の低温段階にPb・Zn等の例がある。金属種を温度だけで一意決定する普遍則ではない。\par

\noindent\textbf{5：代表的条件で適切} 裂け目を通る熱水の鉱脈。流体が運んだ成分が亀裂・断層沿いで沈殿し脈状となる。\par

\paragraph{再利用する考え方} 地質環境・生成機構・物質の名称を対応させ、平均値と全称、典型例と例外、当時の分類と現在の知見を分ける。

\paragraph{共通理論} \hyperref[theory:ore:kuroko]{第\ref*{theory:ore:kuroko}節}

\par 参考：\cite{ore50_kuroko_ref,ore50_skarn_ref,ore50_porphyry_ref}。

\paragraph{関連問題} \hyperref[q:r07-reg-applied-25]{R07-REG-APPLIED-25}

\par 難易度 4（校準初稿）、優先度未評価。

\subsection{令和6年度 問24}
\label{q:r06-reg-applied-24}

\noindent\textbf{原典：} sources/R06/17\_applied\_science.pdf、PDF物理8頁。

\textbf{公式正答：} 3（single）。

\paragraph{問題の要旨} 黒鉱と炭酸塩岩を交代するスカルンを区別する。

\Needspace{4\baselineskip}\paragraph{確認した条件・記述（独自要約）}\mbox{}\par

\noindent\textbf{原典確認：} S7\_050の29原典画像で34問の全五肢・否定・数値を照合。\par

\noindent\textbf{対象：} 黒鉱と炭酸塩岩を交代するスカルンを区別する。\par

\paragraph{問われる核心} 黒鉱と炭酸塩岩を交代するスカルンを区別する。

\paragraph{必要知識} 3は黒鉱をスカルンの生成機構へ置換している。熱水が関与する点は共通でも、海底での硫化物沈殿と、炭酸塩岩を反応させCaケイ酸塩を作る交代は違う。鉱床名・母岩・鉱物・形態・成因を同時に対応させる。

\paragraph{考え方・正答を選ぶ理由（査読待ち初稿）} 3は黒鉱をスカルンの生成機構へ置換している。熱水が関与する点は共通でも、海底での硫化物沈殿と、炭酸塩岩を反応させCaケイ酸塩を作る交代は違う。鉱床名・母岩・鉱物・形態・成因を同時に対応させる。

\Needspace{4\baselineskip}\paragraph{選択肢の検討（査読待ち初稿）}\mbox{}\par

\noindent\textbf{1：代表的条件で適切} 縞状鉄鉱の珪質層と鉄鉱物層。珪質物と鉄鉱物に富む層が交互に現れる。銅の鉱床と交換しない。\par

\noindent\textbf{2：代表的条件で適切} 難溶鉱物の風化残留。風化に伴う移動しにくい成分の相対濃集をいう。\par

\noindent\textbf{3：不適切} 黒鉱が炭酸塩岩の交代で形成。炭酸塩岩との化学反応はスカルンの説明。黒鉱は代表的に海底火山熱水の塊状硫化物である。\par

\noindent\textbf{4：代表的条件で適切} 低品位大規模の斑岩銅。斑岩質貫入体とその周辺の鉱染・網状細脈を持つ代表タイプ。\par

\noindent\textbf{5：代表的条件で適切} 海底熱水からの沈殿。海底火山・熱水活動で溶けた金属が化学条件の変化で沈殿する。\par

\paragraph{再利用する考え方} 地質環境・生成機構・物質の名称を対応させ、平均値と全称、典型例と例外、当時の分類と現在の知見を分ける。

\paragraph{共通理論} \hyperref[theory:ore:kuroko]{第\ref*{theory:ore:kuroko}節}

\par 参考：\cite{ore50_kuroko_ref,ore50_porphyry_ref,ore50_bauxite_ref}。

\paragraph{関連問題} \hyperref[q:r07-reg-applied-25]{R07-REG-APPLIED-25}

\par 難易度 3（校準初稿）、優先度未評価。

\section{天然資源の鉱物・粒径・生成条件を対応させる。}

\subsection{平成25年度 問24}
\label{q:h25-reg-applied-24}

\noindent\textbf{原典：} sources/H25/17\_applied\_science.pdf、PDF物理9頁。

\textbf{公式正答：} 3（single）。

\paragraph{問題の要旨} 天然資源の鉱物・粒径・生成条件を対応させる。

\Needspace{4\baselineskip}\paragraph{確認した条件・記述（独自要約）}\mbox{}\par

\noindent\textbf{原典確認：} S7\_050の29原典画像で34問の全五肢・否定・数値を照合。\par

\noindent\textbf{対象：} 天然資源の鉱物・粒径・生成条件を対応させる。\par

\paragraph{問われる核心} 天然資源の鉱物・粒径・生成条件を対応させる。

\paragraph{必要知識} オイルシェールはケロジェン等を含む細粒堆積岩であり3の粗粒が誤り。縞状鉄鉱、斑岩銅、ガス水和物、風化で濃集したAl鉱石は生成機構が違う。メタンハイドレートのシャーベット状という見た目の表現を、水和物の結晶構造や相安定の定義と同一視しない。

\paragraph{考え方・正答を選ぶ理由（査読待ち初稿）} オイルシェールはケロジェン等を含む細粒堆積岩であり3の粗粒が誤り。縞状鉄鉱、斑岩銅、ガス水和物、風化で濃集したAl鉱石は生成機構が違う。メタンハイドレートのシャーベット状という見た目の表現を、水和物の結晶構造や相安定の定義と同一視しない。

\Needspace{4\baselineskip}\paragraph{選択肢の検討（査読待ち初稿）}\mbox{}\par

\noindent\textbf{1：代表的条件で適切} 縞状鉄鉱の珪質層と鉄鉱物層。珪質物と鉄鉱物に富む層が交互に現れる。銅の鉱床と交換しない。\par

\noindent\textbf{2：代表的条件で適切} 浅所斑岩と斑岩銅。花崗閃緑斑岩等に伴う熱水性の鉱染・網状細脈を作る。\par

\noindent\textbf{3：不適切} 粗粒堆積岩がオイルシェール。ケロジェンを含む細粒堆積岩として扱う。粗粒という説明が不適切。\par

\noindent\textbf{4：代表的条件で適切} メタンハイドレートの温度圧力条件。水分子の構造にガス分子が取り込まれた水和物で、低温・高圧等の安定条件がある。単に凍結したメタンではない。\par

\noindent\textbf{5：代表的条件で適切} Al水酸化物がボーキサイト。Alの水酸化物・酸化水酸化物等を含む主要なAl鉱石。原岩が風化してAlが相対濃集する。\par

\paragraph{再利用する考え方} 地質環境・生成機構・物質の名称を対応させ、平均値と全称、典型例と例外、当時の分類と現在の知見を分ける。

\paragraph{共通理論} \hyperref[theory:ore:sedimentary]{第\ref*{theory:ore:sedimentary}節}

\par 参考：\cite{ore50_oilshale_ref,ore50_bauxite_ref,ore50_bif_ref}。

\paragraph{関連問題} \hyperref[q:h28-reg-applied-31]{H28-REG-APPLIED-31}

\par 難易度 3（校準初稿）、優先度未評価。

\subsection{平成28年度 問31}
\label{q:h28-reg-applied-31}

\noindent\textbf{原典：} sources/H28/17\_applied\_science.pdf、PDF物理12頁。

\textbf{公式正答：} 3（single）。

\paragraph{問題の要旨} 天然資源の鉱物・粒径・生成条件を対応させる。

\Needspace{4\baselineskip}\paragraph{確認した条件・記述（独自要約）}\mbox{}\par

\noindent\textbf{原典確認：} S7\_050の29原典画像で34問の全五肢・否定・数値を照合。\par

\noindent\textbf{対象：} 天然資源の鉱物・粒径・生成条件を対応させる。\par

\paragraph{問われる核心} 天然資源の鉱物・粒径・生成条件を対応させる。

\paragraph{必要知識} オイルシェールはケロジェン等を含む細粒堆積岩であり3の粗粒が誤り。縞状鉄鉱、斑岩銅、ガス水和物、風化で濃集したAl鉱石は生成機構が違う。メタンハイドレートのシャーベット状という見た目の表現を、水和物の結晶構造や相安定の定義と同一視しない。

\paragraph{考え方・正答を選ぶ理由（査読待ち初稿）} オイルシェールはケロジェン等を含む細粒堆積岩であり3の粗粒が誤り。縞状鉄鉱、斑岩銅、ガス水和物、風化で濃集したAl鉱石は生成機構が違う。メタンハイドレートのシャーベット状という見た目の表現を、水和物の結晶構造や相安定の定義と同一視しない。

\Needspace{4\baselineskip}\paragraph{選択肢の検討（査読待ち初稿）}\mbox{}\par

\noindent\textbf{1：代表的条件で適切} 縞状鉄鉱の珪質層と鉄鉱物層。珪質物と鉄鉱物に富む層が交互に現れる。銅の鉱床と交換しない。\par

\noindent\textbf{2：代表的条件で適切} 浅所斑岩と斑岩銅。花崗閃緑斑岩等に伴う熱水性の鉱染・網状細脈を作る。\par

\noindent\textbf{3：不適切} 粗粒堆積岩がオイルシェール。ケロジェンを含む細粒堆積岩として扱う。粗粒という説明が不適切。\par

\noindent\textbf{4：代表的条件で適切} メタンハイドレートの温度圧力条件。水分子の構造にガス分子が取り込まれた水和物で、低温・高圧等の安定条件がある。単に凍結したメタンではない。\par

\noindent\textbf{5：代表的条件で適切} Al水酸化物がボーキサイト。Alの水酸化物・酸化水酸化物等を含む主要なAl鉱石。原岩が風化してAlが相対濃集する。\par

\paragraph{再利用する考え方} 地質環境・生成機構・物質の名称を対応させ、平均値と全称、典型例と例外、当時の分類と現在の知見を分ける。

\paragraph{共通理論} \hyperref[theory:ore:sedimentary]{第\ref*{theory:ore:sedimentary}節}

\par 参考：\cite{ore50_oilshale_ref,ore50_bauxite_ref,ore50_bif_ref}。

\paragraph{関連問題} \hyperref[q:h25-reg-applied-24]{H25-REG-APPLIED-24}

\par 難易度 3（校準初稿）、優先度未評価。

\section{マンガン団塊の組成・海底分布・成長時間を比較する。}

\subsection{平成23年度 問27}
\label{q:h23-reg-applied-27}

\noindent\textbf{原典：} sources/H23/17\_applied\_science.pdf、PDF物理9頁。

\textbf{公式正答：} 5（single）。

\par\textbf{検討未確定：} 副成分の最大1–2\%という上限と粒径の全分類範囲は未確定。成長も海水起源・続成作用等で異なり、一つの代表率を全試料へ固定しない。

\paragraph{問題の要旨} マンガン団塊の組成・海底分布・成長時間を比較する。

\Needspace{4\baselineskip}\paragraph{確認した条件・記述（独自要約）}\mbox{}\par

\noindent\textbf{原典確認：} S7\_050の29原典画像で34問の全五肢・否定・数値を照合。\par

\noindent\textbf{対象：} マンガン団塊の組成・海底分布・成長時間を比較する。\par

\paragraph{問われる核心} マンガン団塊の組成・海底分布・成長時間を比較する。

\paragraph{必要知識} 団塊は核の周囲へ鉄・マンガン酸化物等がゆっくり堆積する。JOGMECの例は百万年で1mm程度で、5の数mm／1万年は時間尺度が大きく違う。海底に露出した団塊と岩盤上のクラストを区別し、成長率は起源・環境で変わる代表値として読む。

\paragraph{考え方・正答を選ぶ理由（査読待ち初稿）} 団塊は核の周囲へ鉄・マンガン酸化物等がゆっくり堆積する。JOGMECの例は百万年で1mm程度で、5の数mm／1万年は時間尺度が大きく違う。海底に露出した団塊と岩盤上のクラストを区別し、成長率は起源・環境で変わる代表値として読む。

\Needspace{4\baselineskip}\paragraph{選択肢の検討（査読待ち初稿）}\mbox{}\par

\noindent\textbf{1：代表的条件で適切} Mn・Fe酸化物の黒色団塊。鉄・マンガン酸化物を主体に成長する団塊。原典の1mm以上は分類上の目安として保持する。\par

\noindent\textbf{2：条件・解釈を留保} Cu・Ni・Coを1–2\%含む場合。これら金属を副成分として含む。原典の最大1–2\%を全海域・全団塊の厳密上限とする根拠は未確認。\par

\noindent\textbf{3：代表的条件で適切} 堆積の遅い深海底の表層。団塊の埋没が遅い場所で海水・間隙水から成長する代表的な産状。露岩を覆うクラストと区別する。\par

\noindent\textbf{4：代表的条件で適切} 岩片や化石を核に成長。核を中心に層状に成長する場合がある。\par

\noindent\textbf{5：不適切} 数mm／1万年の成長。代表値は百万年で1mm程度と極めて遅い。原典の万年は代表的な海水沈殿成長の時間尺度と合わない。\par

\paragraph{再利用する考え方} 地質環境・生成機構・物質の名称を対応させ、平均値と全称、典型例と例外、当時の分類と現在の知見を分ける。

\paragraph{共通理論} \hyperref[theory:ore:marine]{第\ref*{theory:ore:marine}節}

\par 参考：\cite{ore50_nodule_ref}。

\par 難易度 4（校準初稿）、優先度未評価。

\section{風化・堆積・マグマ・熱水・石油の濃集機構を比較する。}

\subsection{平成27年度 問31}
\label{q:h27-reg-applied-31}

\noindent\textbf{原典：} sources/H27/17\_applied\_science.pdf、PDF物理12頁。

\textbf{公式正答：} 1（single）。

\paragraph{問題の要旨} 風化・堆積・マグマ・熱水・石油の濃集機構を比較する。

\Needspace{4\baselineskip}\paragraph{確認した条件・記述（独自要約）}\mbox{}\par

\noindent\textbf{原典確認：} S7\_050の29原典画像で34問の全五肢・否定・数値を照合。\par

\noindent\textbf{対象：} 風化・堆積・マグマ・熱水・石油の濃集機構を比較する。\par

\paragraph{問われる核心} 風化・堆積・マグマ・熱水・石油の濃集機構を比較する。

\paragraph{必要知識} 岩石から何が移動し、何が残るかを追う。風化残留でAl・Feが全て溶出し、Na・K・Ca・Mgが残るという交換が誤りである。他の成因は、粒子や溶存成分の堆積、硫化物液の分離と結晶集積、熱水の沈殿、炭化水素の移動とトラップとして区別する。pH・酸化還元条件で元素の移動性が変わる例まで固定した順位にしない。

\paragraph{考え方・正答を選ぶ理由（査読待ち初稿）} 岩石から何が移動し、何が残るかを追う。風化残留でAl・Feが全て溶出し、Na・K・Ca・Mgが残るという交換が誤りである。他の成因は、粒子や溶存成分の堆積、硫化物液の分離と結晶集積、熱水の沈殿、炭化水素の移動とトラップとして区別する。pH・酸化還元条件で元素の移動性が変わる例まで固定した順位にしない。

\Needspace{4\baselineskip}\paragraph{選択肢の検討（査読待ち初稿）}\mbox{}\par

\noindent\textbf{1：不適切} Al・Feを溶かしてNa・K・Ca・Mgを残留。通常の強い風化残留では移動しやすい成分が溶脱し、Al・Fe等の難溶性成分が相対濃集する。列挙した方向が逆。\par

\noindent\textbf{2：代表的条件で適切} 溶存成分沈殿又は粒子堆積。化学的な沈殿と物理的な運搬・分級による堆積鉱床の二つを含む。\par

\noindent\textbf{3：代表的条件で適切} 硫化物メルト分離・結晶集積。マグマ内の液相不混和や結晶分化・集積で金属が濃集する正マグマ成因。\par

\noindent\textbf{4：代表的条件で適切} 熱水中の元素の濃縮・沈殿。温度・圧力・酸化還元・pH等の変化で溶解度が変わり鉱化する。\par

\noindent\textbf{5：代表的条件で適切} 石油のトラップ濃集。移動する炭化水素が貯留層と封じる構造へ集まる。根源岩での生成とトラップは別段階。\par

\paragraph{再利用する考え方} 地質環境・生成機構・物質の名称を対応させ、平均値と全称、典型例と例外、当時の分類と現在の知見を分ける。

\paragraph{共通理論} \hyperref[theory:ore:origins]{第\ref*{theory:ore:origins}節}

\par 参考：\cite{ore50_bauxite_ref,ore50_porphyry_ref,ore50_magmatic_ref,ore50_besshi_ref}。

\paragraph{関連問題} \hyperref[q:r02-reg-applied-26]{R02-REG-APPLIED-26}

\par 難易度 3（校準初稿）、優先度未評価。

\subsection{平成29年度 問28}
\label{q:h29-reg-applied-28}

\noindent\textbf{原典：} sources/H29/17\_applied\_science.pdf、PDF物理9頁。

\textbf{公式正答：} 1（single）。

\par\textbf{検討未確定：} 公式1を保持。原典3のウランを一般の現代河川での漂砂と読むと成立しにくい。古代石英礫岩の還元環境と砂岩型の化学沈殿を区別する必要がある。 原典1は金銀の早期結晶がある場合と述べるため、代表例がCr/Ptであるだけでは全可能性を排除できない。マグマ性PGE鉱床に自然金を含むモデルもあり、公式1の排他理由を要査読。

\paragraph{問題の要旨} 正マグマ濃集・漂砂・熱水・風化・交代を比較する。

\Needspace{4\baselineskip}\paragraph{確認した条件・記述（独自要約）}\mbox{}\par

\noindent\textbf{原典確認：} S7\_050の29原典画像で34問の全五肢・否定・数値を照合。\par

\noindent\textbf{対象：} 正マグマ濃集・漂砂・熱水・風化・交代を比較する。\par

\paragraph{問われる核心} 正マグマ濃集・漂砂・熱水・風化・交代を比較する。

\paragraph{必要知識} 生成機構は正マグマの結晶・硫化物液集積、熱水の沈殿、粒子の漂砂、風化残留、母岩の化学的交代に分ける。早期に金銀を集積するという1を、クロム・白金等の代表例と区別して選ぶ。ウランの漂砂は酸化されにくい古代の環境等を条件とし、現代の河川で全ウラン鉱が砂金と同様に残ると説明しない。

\paragraph{考え方・正答を選ぶ理由（査読待ち初稿）} 生成機構は正マグマの結晶・硫化物液集積、熱水の沈殿、粒子の漂砂、風化残留、母岩の化学的交代に分ける。早期に金銀を集積するという1を、クロム・白金等の代表例と区別して選ぶ。ウランの漂砂は酸化されにくい古代の環境等を条件とし、現代の河川で全ウラン鉱が砂金と同様に残ると説明しない。

\Needspace{4\baselineskip}\paragraph{選択肢の検討（査読待ち初稿）}\mbox{}\par

\noindent\textbf{1：公式の選択・一般化を留保} 早期結晶する金・銀の正マグマ鉱。早期集積の代表はクロム・白金等。金銀鉱床は熱水等の成因が重要で、原典の例示が適切でない。 ただし原典は場合があるという存在命題で、全ての金銀マグマ成因を否定する論証にはならない。\par

\noindent\textbf{2：代表的条件で適切} 金属を運ぶ熱水。亀裂中又は海底等で金属成分を沈殿させる。全熱水鉱床の流体をマグマ起源だけに限定しない。\par

\noindent\textbf{3：条件・解釈を留保} ウラン鉱と砂金の漂砂濃集。機械的な粒子の運搬・堆積による砂金は代表例。ウラニナイトの漂砂は酸素の少ない古い時代等の条件が必要で、砂岩中の溶存Uの酸化還元沈殿とは違う。\par

\noindent\textbf{4：代表的条件で適切} 難溶物の風化残留。母岩の風化で移動しにくい成分が相対濃集し、ボーキサイト等となる。\par

\noindent\textbf{5：代表的条件で適切} 石灰岩との接触交代。熱水等との反応でCaを含むケイ酸塩鉱物を作るスカルン成因。単なる再結晶と分ける。\par

\paragraph{再利用する考え方} 地質環境・生成機構・物質の名称を対応させ、平均値と全称、典型例と例外、当時の分類と現在の知見を分ける。

\paragraph{共通理論} \hyperref[theory:ore:origins]{第\ref*{theory:ore:origins}節}

\par 参考：\cite{ore50_porphyry_ref,ore50_bauxite_ref,ore50_uranium_ref,ore50_magmatic_ref,ore50_besshi_ref,ore50_gold_exception_ref}。

\par 難易度 4（校準初稿）、優先度未評価。

\subsection{令和2年度 問26}
\label{q:r02-reg-applied-26}

\noindent\textbf{原典：} sources/R02/17\_applied\_science.pdf、PDF物理10頁。

\textbf{公式正答：} 3（single）。

\paragraph{問題の要旨} 風化・堆積・マグマ・熱水・石油の濃集機構を比較する。

\Needspace{4\baselineskip}\paragraph{確認した条件・記述（独自要約）}\mbox{}\par

\noindent\textbf{原典確認：} S7\_050の29原典画像で34問の全五肢・否定・数値を照合。\par

\noindent\textbf{対象：} 風化・堆積・マグマ・熱水・石油の濃集機構を比較する。\par

\paragraph{問われる核心} 風化・堆積・マグマ・熱水・石油の濃集機構を比較する。

\paragraph{必要知識} 岩石から何が移動し、何が残るかを追う。風化残留でAl・Feが全て溶出し、Na・K・Ca・Mgが残るという交換が誤りである。他の成因は、粒子や溶存成分の堆積、硫化物液の分離と結晶集積、熱水の沈殿、炭化水素の移動とトラップとして区別する。pH・酸化還元条件で元素の移動性が変わる例まで固定した順位にしない。

\paragraph{考え方・正答を選ぶ理由（査読待ち初稿）} 岩石から何が移動し、何が残るかを追う。風化残留でAl・Feが全て溶出し、Na・K・Ca・Mgが残るという交換が誤りである。他の成因は、粒子や溶存成分の堆積、硫化物液の分離と結晶集積、熱水の沈殿、炭化水素の移動とトラップとして区別する。pH・酸化還元条件で元素の移動性が変わる例まで固定した順位にしない。

\Needspace{4\baselineskip}\paragraph{選択肢の検討（査読待ち初稿）}\mbox{}\par

\noindent\textbf{1：代表的条件で適切} 熱水中の元素の濃縮・沈殿。温度・圧力・酸化還元・pH等の変化で溶解度が変わり鉱化する。\par

\noindent\textbf{2：代表的条件で適切} 硫化物メルト分離・結晶集積。マグマ内の液相不混和や結晶分化・集積で金属が濃集する正マグマ成因。\par

\noindent\textbf{3：不適切} Al・Feを溶かしてNa・K・Ca・Mgを残留。通常の強い風化残留では移動しやすい成分が溶脱し、Al・Fe等の難溶性成分が相対濃集する。列挙した方向が逆。\par

\noindent\textbf{4：代表的条件で適切} 石油のトラップ濃集。移動する炭化水素が貯留層と封じる構造へ集まる。根源岩での生成とトラップは別段階。\par

\noindent\textbf{5：代表的条件で適切} 溶存成分沈殿又は粒子堆積。化学的な沈殿と物理的な運搬・分級による堆積鉱床の二つを含む。\par

\paragraph{再利用する考え方} 地質環境・生成機構・物質の名称を対応させ、平均値と全称、典型例と例外、当時の分類と現在の知見を分ける。

\paragraph{共通理論} \hyperref[theory:ore:origins]{第\ref*{theory:ore:origins}節}

\par 参考：\cite{ore50_bauxite_ref,ore50_porphyry_ref,ore50_magmatic_ref,ore50_besshi_ref}。

\paragraph{関連問題} \hyperref[q:h27-reg-applied-31]{H27-REG-APPLIED-31}

\par 難易度 3（校準初稿）、優先度未評価。

\subsection{令和7年度 問25}
\label{q:r07-reg-applied-25}

\noindent\textbf{原典：} sources/R07/17\_applied\_science.pdf、PDF物理9頁。

\textbf{公式正答：} 3（single）。

\paragraph{問題の要旨} 正マグマ・黒鉱・別子型・スカルンの産状を比較する。

\Needspace{4\baselineskip}\paragraph{確認した条件・記述（独自要約）}\mbox{}\par

\noindent\textbf{原典確認：} S7\_050の29原典画像で34問の全五肢・否定・数値を照合。\par

\noindent\textbf{対象：} 正マグマ・黒鉱・別子型・スカルンの産状を比較する。\par

\paragraph{問われる核心} 正マグマ・黒鉱・別子型・スカルンの産状を比較する。

\paragraph{必要知識} 原典3の黒鉱の母岩が誤りで、緑色片岩ではなく代表的には中新世グリーンタフ地域である。2は正マグマ、4は変成を経た海底硫化物、5は炭酸塩岩の化学交代であり、同じ金属資源でも生成機構を別々に読む。

\paragraph{考え方・正答を選ぶ理由（査読待ち初稿）} 原典3の黒鉱の母岩が誤りで、緑色片岩ではなく代表的には中新世グリーンタフ地域である。2は正マグマ、4は変成を経た海底硫化物、5は炭酸塩岩の化学交代であり、同じ金属資源でも生成機構を別々に読む。

\Needspace{4\baselineskip}\paragraph{選択肢の検討（査読待ち初稿）}\mbox{}\par

\noindent\textbf{1：代表的条件で適切} 通常より鉱物・元素が濃集した鉱床。原典では地質学上の濃集をいう。経済的に採掘できるかは品位・鉱量・価格・技術等を別に評価する。\par

\noindent\textbf{2：代表的条件で適切} Cr・Ptの早期結晶と集積。苦鉄質～超苦鉄質のマグマ系でクロム・白金族等が濃集する代表的な正マグマ鉱床。全鉱体が厳密に下部だけという意味ではない。\par

\noindent\textbf{3：不適切} 黒鉱を東北の緑色片岩中とする。グリーンタフ地域の海底火山熱水に伴う黒鉱と、緑色片岩という変成岩を交換している。\par

\noindent\textbf{4：代表的条件で適切} 海底熱水鉱が変成した層状含銅硫化鉄鉱。海洋地殻上の熱水性硫化物が付加・沈み込みと変成を経る別子型等の代表モデル。\par

\noindent\textbf{5：代表的条件で適切} 炭酸塩との反応でCaケイ酸塩を作るスカルン。熱水による化学的交代で形成される。単なる熱だけの再結晶と分ける。\par

\paragraph{再利用する考え方} 地質環境・生成機構・物質の名称を対応させ、平均値と全称、典型例と例外、当時の分類と現在の知見を分ける。

\paragraph{共通理論} \hyperref[theory:ore:origins]{第\ref*{theory:ore:origins}節}

\par 参考：\cite{ore50_kuroko_ref,ore50_porphyry_ref,ore50_skarn_ref,ore50_magmatic_ref,ore50_besshi_ref,ore50_besshi_ref}。

\paragraph{関連問題} \hyperref[q:r05-reg-applied-23]{R05-REG-APPLIED-23}, \hyperref[q:r06-reg-applied-24]{R06-REG-APPLIED-24}

\par 難易度 3（校準初稿）、優先度未評価。

\section{斑岩銅・黒鉱・別子型・スカルンの生成場を比較する。}

\subsection{平成23年度 問23}
\label{q:h23-reg-applied-23}

\noindent\textbf{原典：} sources/H23/17\_applied\_science.pdf、PDF物理7頁。

\textbf{公式正答：} 1（single）。

\paragraph{問題の要旨} 斑岩銅・黒鉱・別子型・スカルンの生成場を比較する。

\Needspace{4\baselineskip}\paragraph{確認した条件・記述（独自要約）}\mbox{}\par

\noindent\textbf{原典確認：} S7\_050の29原典画像で34問の全五肢・否定・数値を照合。\par

\noindent\textbf{対象：} 斑岩銅・黒鉱・別子型・スカルンの生成場を比較する。\par

\paragraph{問われる核心} 斑岩銅・黒鉱・別子型・スカルンの生成場を比較する。

\paragraph{必要知識} 斑岩銅は浅所の斑岩質貫入と熱水系に関係し、岩体内外の鉱染・網状細脈として形成されるので1。硫化物を含むという共通点だけで、海底の塊状鉱、変成した層状鉱、炭酸塩岩を交代したスカルンを統合しない。

\paragraph{考え方・正答を選ぶ理由（査読待ち初稿）} 斑岩銅は浅所の斑岩質貫入と熱水系に関係し、岩体内外の鉱染・網状細脈として形成されるので1。硫化物を含むという共通点だけで、海底の塊状鉱、変成した層状鉱、炭酸塩岩を交代したスカルンを統合しない。

\Needspace{4\baselineskip}\paragraph{選択肢の検討（査読待ち初稿）}\mbox{}\par

\noindent\textbf{1：代表的条件で適切} 浅所貫入岩内外の鉱染・細脈。斑岩質岩体とその周辺に分散した硫化物と網状細脈を作り、大規模・低品位となる代表的な斑岩銅の産状。\par

\noindent\textbf{2：不適切} 珪長質海底火山の塊状硫化物。火山性塊状硫化物鉱床、特に黒鉱型の説明。斑岩銅と生成位置・形態が異なる。\par

\noindent\textbf{3：不適切} 苦鉄質起源の変成岩中の層状硫化物。別子型等の層状含銅硫化鉄鉱の説明に対応する。\par

\noindent\textbf{4：不適切} 炭酸塩岩の熱水交代。接触交代・スカルン鉱床の説明で、斑岩銅の鉱染・細脈とは区別する。\par

\noindent\textbf{5：不適切} 縞状鉱でCu2+を酸化沈殿。古い縞状鉄鉱の成因を銅へ置換している。鉄の酸化・沈殿を述べる場合と混同しない。\par

\paragraph{再利用する考え方} 地質環境・生成機構・物質の名称を対応させ、平均値と全称、典型例と例外、当時の分類と現在の知見を分ける。

\paragraph{共通理論} \hyperref[theory:ore:hydrothermal]{第\ref*{theory:ore:hydrothermal}節}

\par 参考：\cite{ore50_porphyry_ref,ore50_magmatic_ref,ore50_besshi_ref}。

\par 難易度 3（校準初稿）、優先度未評価。

\section{収束・発散境界と黒鉱・斑岩銅の関係を照合する。}

\subsection{平成24年度 問25}
\label{q:h24-reg-applied-25}

\noindent\textbf{原典：} sources/H24/17\_applied\_science.pdf、PDF物理8頁。

\textbf{公式正答：} 4（single）。

\par\textbf{独立解答と公式正答に相違あり：} 独立解答 2。初回2／公式4。正答原典p41画像も確認。海洋性島弧の斑岩銅と背弧海底の黒鉱型があり、肢2と4の排他性及びマリアナ型の当時定義は要査読。初回を公式へ合わせて置換しない。

\paragraph{問題の要旨} 収束・発散境界と黒鉱・斑岩銅の関係を照合する。

\Needspace{4\baselineskip}\paragraph{確認した条件・記述（独自要約）}\mbox{}\par

\noindent\textbf{原典確認：} S7\_050の29原典画像で34問の全五肢・否定・数値を照合。\par

\noindent\textbf{対象：} 収束・発散境界と黒鉱・斑岩銅の関係を照合する。\par

\paragraph{問われる核心} 収束・発散境界と黒鉱・斑岩銅の関係を照合する。

\paragraph{必要知識} 1と3は収束・発散を交換している。5は堆積盆地の砂岩型ウランをプレート境界へ一対一対応させている。公式原典の平成24年度正答表41頁で4を再確認した。一方、黒鉱型の海底熱水活動は島弧・背弧にも見られ、初回独立2を排除する根拠が不足している。斑岩銅をアンデス型だけと断定することも避け、マリアナ型という用語と鉱床タイプの対応を査読に残す。

\paragraph{考え方・正答を選ぶ理由（査読待ち初稿）} 1と3は収束・発散を交換している。5は堆積盆地の砂岩型ウランをプレート境界へ一対一対応させている。公式原典の平成24年度正答表41頁で4を再確認した。一方、黒鉱型の海底熱水活動は島弧・背弧にも見られ、初回独立2を排除する根拠が不足している。斑岩銅をアンデス型だけと断定することも避け、マリアナ型という用語と鉱床タイプの対応を査読に残す。

\Needspace{4\baselineskip}\paragraph{選択肢の検討（査読待ち初稿）}\mbox{}\par

\noindent\textbf{1：不適切} 発散境界の島弧海溝。島弧海溝系はプレート収束と沈み込みに関係する。発散は海嶺等の説明。\par

\noindent\textbf{2：条件・解釈を留保} マリアナ型と黒鉱。海洋性島弧や伸張する背弧の海底熱水活動に黒鉱型の潜在鉱床がある。マリアナ型の定義を固定しない限りこの肢を簡単に排除できない。\par

\noindent\textbf{3：不適切} 収束境界の中央海嶺。新しい海洋地殻を作る中央海嶺は発散境界である。\par

\noindent\textbf{4：条件・解釈を留保} マリアナ型と斑岩銅。斑岩銅は収束縁の島弧・陸弧等に産する。公式は4だが、マリアナ型への特定と2との差を一次資料から一意に説明する確認は未完了。\par

\noindent\textbf{5：不適切} 収束境界に砂岩型ウラン。砂岩中の流体移動と酸化還元境界等を使う分類であり、収束プレート境界だけを成立条件としない。\par

\paragraph{再利用する考え方} 地質環境・生成機構・物質の名称を対応させ、平均値と全称、典型例と例外、当時の分類と現在の知見を分ける。

\paragraph{共通理論} \hyperref[theory:ore:tectonic]{第\ref*{theory:ore:tectonic}節}

\par 参考：\cite{ore50_porphyry_ref,ore50_backarc_ref,usgs_plate_boundaries}。

\par 難易度 4（校準初稿）、優先度未評価。

\section{ウラン鉱床の母岩・年代・品位の分類を比較する。}

\subsection{令和4年度 問24}
\label{q:r04-reg-applied-24}

\noindent\textbf{原典：} sources/R04/17\_applied\_science.pdf、PDF物理8頁。

\textbf{公式正答：} 5（single）。

\par\textbf{検討未確定：} 公式5を保持。4億年以降と23億年前以前等の厳密な分類時代境界・例外は原典の当時分類とIAEA原分類表の追加照合が必要。低品位資源の存在のみで全鉱床の一律平均品位を作らない。

\paragraph{問題の要旨} ウラン鉱床の母岩・年代・品位の分類を比較する。

\Needspace{4\baselineskip}\paragraph{確認した条件・記述（独自要約）}\mbox{}\par

\noindent\textbf{原典確認：} S7\_050の29原典画像で34問の全五肢・否定・数値を照合。\par

\noindent\textbf{対象：} ウラン鉱床の母岩・年代・品位の分類を比較する。\par

\paragraph{問われる核心} ウラン鉱床の母岩・年代・品位の分類を比較する。

\paragraph{必要知識} ウラン濃集は母岩・流体・酸化還元と形態で分類する。貫入岩型には低～非常に低品位の大規模資源があり、全て高品位とした5を選ぶ。古い石英礫岩の漂砂、不整合近傍の流体、砂岩の還元帯、鉱脈の亀裂充填を区別する。品位と鉱量は別の評価量である。

\paragraph{考え方・正答を選ぶ理由（査読待ち初稿）} ウラン濃集は母岩・流体・酸化還元と形態で分類する。貫入岩型には低～非常に低品位の大規模資源があり、全て高品位とした5を選ぶ。古い石英礫岩の漂砂、不整合近傍の流体、砂岩の還元帯、鉱脈の亀裂充填を区別する。品位と鉱量は別の評価量である。

\Needspace{4\baselineskip}\paragraph{選択肢の検討（査読待ち初稿）}\mbox{}\par

\noindent\textbf{1：代表的条件で適切} 古い基盤と砂岩の不整合近傍。不整合関連型の代表は原生代砂岩盆地と古い変成基盤の不整合をまたぐ熱水・流体鉱化。\par

\noindent\textbf{2：条件・解釈を留保} 4億年前以降の河川砂岩・礫岩。顕生代の堆積盆地は砂岩型の代表だが、4億年を全砂岩型鉱床の唯一の年齢限界とする確認は未完了。\par

\noindent\textbf{3：条件・解釈を留保} 約23億年前以前の石英礫岩。古い酸素の少ない環境でU鉱物等が漂砂濃集する代表的な古代石英礫岩型。年代は酸化環境と地域史の条件として扱う。\par

\noindent\textbf{4：代表的条件で適切} 様々な時代の鉱脈型。亀裂等を切る流体の通路に脈状鉱体ができる形態で、特定時代だけに限定されない。\par

\noindent\textbf{5：不適切} 貫入岩型は他型より高品位。一般に低品位で大規模な貫入岩関連資源があり、高品位を共通特徴としない。\par

\paragraph{再利用する考え方} 地質環境・生成機構・物質の名称を対応させ、平均値と全称、典型例と例外、当時の分類と現在の知見を分ける。

\paragraph{共通理論} \hyperref[theory:ore:uranium]{第\ref*{theory:ore:uranium}節}

\par 参考：\cite{ore50_uranium_ref,ore50_intrusive_ref}。

\par 難易度 4（校準初稿）、優先度未評価。
